% =====================================================================
%  courses/trigonometria/indice.tex — índice detallado: trigonometría preuniversitaria
% ---------------------------------------------------------------------
%  5 Unidades · 21 Temas
% =====================================================================

{
\hypersetup{linkcolor=AzulPizarra}
\tableofcontents
}
\newpage

\thispagestyle{fancy}

\begin{center}
  {\Large\sffamily\bfseries\color{AzulPizarra} Índice Temático — Trigonometría Preuniversitaria}\\[4pt]
  {\small\sffamily\color{GrisMedio} 5 Unidades · 21 Temas · Teoría, Ejercicios y Resueltos}
\end{center}

\vspace{0.6cm}

% --- UNIDAD I --------------------------------------------------------------
\begin{tcolorbox}[enhanced, colback=AzulPizarra, colframe=DoradoAcadémico,
    arc=3pt, boxrule=1pt, left=8pt, right=8pt, top=5pt, bottom=5pt,
    fontupper=\sffamily\bfseries\color{white}]
  UNIDAD I \hfill Fundamentos \hfill Temas 1–4
\end{tcolorbox}
\begin{multicols}{2}
  \begin{enumerate}[label=\textcolor{AzulPizarra}{\textbf{\arabic*.}},
      leftmargin=2em, itemsep=1pt]
    \item Ángulos
    \item Sistemas de medición angular
    \item Conversión de unidades
    \item Arcos y sectores
  \end{enumerate}
\end{multicols}

\vspace{0.4cm}

% --- UNIDAD II -------------------------------------------------------------
\begin{tcolorbox}[enhanced, colback=AzulMedio, colframe=DoradoAcadémico,
    arc=3pt, boxrule=1pt, left=8pt, right=8pt, top=5pt, bottom=5pt,
    fontupper=\sffamily\bfseries\color{white}]
  UNIDAD II \hfill Razones Trigonométricas \hfill Temas 5–8
\end{tcolorbox}
\begin{multicols}{2}
  \begin{enumerate}[label=\textcolor{AzulMedio}{\textbf{\arabic*.}},
      leftmargin=2em, itemsep=1pt, start=5]
    \item Razones trigonométricas
    \item Razones en triángulos rectángulos
    \item Razones de ángulos notables
    \item Razones trigonométricas recíprocas
  \end{enumerate}
\end{multicols}

\vspace{0.4cm}

% --- UNIDAD III ------------------------------------------------------------
\begin{tcolorbox}[enhanced, colback=AzulPizarra, colframe=DoradoAcadémico,
    arc=3pt, boxrule=1pt, left=8pt, right=8pt, top=5pt, bottom=5pt,
    fontupper=\sffamily\bfseries\color{white}]
  UNIDAD III \hfill Identidades \hfill Temas 9–12
\end{tcolorbox}
\begin{multicols}{2}
  \begin{enumerate}[label=\textcolor{AzulPizarra}{\textbf{\arabic*.}},
      leftmargin=2em, itemsep=1pt, start=9]
    \item Identidades trigonométricas
    \item Identidades fundamentales
    \item Identidades auxiliares
    \item Transformaciones trigonométricas
  \end{enumerate}
\end{multicols}

\vspace{0.4cm}

% --- UNIDAD IV -------------------------------------------------------------
\begin{tcolorbox}[enhanced, colback=AzulMedio, colframe=DoradoAcadémico,
    arc=3pt, boxrule=1pt, left=8pt, right=8pt, top=5pt, bottom=5pt,
    fontupper=\sffamily\bfseries\color{white}]
  UNIDAD IV \hfill Ecuaciones y Funciones \hfill Temas 13–16
\end{tcolorbox}
\begin{multicols}{2}
  \begin{enumerate}[label=\textcolor{AzulMedio}{\textbf{\arabic*.}},
      leftmargin=2em, itemsep=1pt, start=13]
    \item Ecuaciones trigonométricas
    \item Inecuaciones trigonométricas
    \item Funciones trigonométricas
    \item Gráficas trigonométricas
  \end{enumerate}
\end{multicols}

\vspace{0.4cm}

% --- UNIDAD V --------------------------------------------------------------
\begin{tcolorbox}[enhanced, colback=AzulPizarra, colframe=DoradoAcadémico,
    arc=3pt, boxrule=1pt, left=8pt, right=8pt, top=5pt, bottom=5pt,
    fontupper=\sffamily\bfseries\color{white}]
  UNIDAD V \hfill Trigonometría Analítica \hfill Temas 17–21
\end{tcolorbox}
\begin{multicols}{2}
  \begin{enumerate}[label=\textcolor{AzulPizarra}{\textbf{\arabic*.}},
      leftmargin=2em, itemsep=1pt, start=17]
    \item Ley de senos
    \item Ley de cosenos
    \item Resolución de triángulos
    \item Área de triángulos
    \item Fórmulas trigonométricas avanzadas
  \end{enumerate}
\end{multicols}

\vspace{0.6cm}

\begin{cajatip}
  La trigonometría es la puerta de entrada a la física ondulatoria, la geometría
  analítica y el cálculo. Dominar las identidades fundamentales (Unidad~III)
  permite resolver el 70\% de los problemas de examen de admisión sin necesidad
  de memorizar fórmulas adicionales.
\end{cajatip}

\newpage